% Part: proof-theory
% Chapter: normalization
% Section: reductions

\documentclass[../../../include/open-logic-section]{subfiles}

\begin{document}

\olfileid{pt}{nor}{red}

\olsection{د کمولو بدلونونه}

لومړے د معرفي کولو د قاعدې حالت اخلو۔ که په \Log{N1i} کښې د
!!a{proof} د پرېکون اوږدوالے~$1$ وي، دا يوه داسې برخه ده چې د
!!{formula} يو مورد لري، او همدغه !!{formula} هم د
!!a{introduction} استنتاج پايله او هم د !!a{elimination} استنتاج
لويه مقدمه ده۔ د عادي کولو په ثبوت کښې دا ډول پرېکونونه داسې
``کمول'' کېږي چې په دغې !!{introduction}/!!{elimination} جوړې
پاے ته رسېدونکے فرعي !!{proof} په نوي !!{proof} بدل شي، چې
د استنتاجونو دا جوړه ترې لرې شوې وي۔ داسې بدلول نوي پرېکونونه
پيدا کولے شي۔ خو کتلے شو چې د نويو پرېکونونو رتبه د کمېدونکي
پرېکون له رتبې ټيټه ده۔ نو د نوي !!{proof} يا د پرېکون اوږدوالے
کم دے، ځکه د تر ټولو لوړې رتبې يو د~$1$ اوږدوالي پرېکون لرې شو،
يا د !!{proof} د پرېکون رتبه ټيټه شوې ده، که همدا د تر ټولو
لوړې رتبې يوازينے پرېکون و۔

مثلاً، که اړوند استنتاجونه \Intro{\lif} او \Elim{\lif}
وي، $!A \ident !B \lif !C$ وي او داسې فرعي
!!{proof}~$\delta_1$ وي چې په لاندې رسم پاے ته رسېږي:
\[
\AxiomC{$\Discharge{!B}{x}$}
\RightLabel{$\delta_2$}
\DeduceC{$!C$}
\DischargeRule{\Intro{\lif}}{x}
\UnaryInfC{$!B \lif !C$}
\AxiomC{}
\RightLabel{$\delta_3$}
\DeduceC{$!B$}
\RightLabel{\Elim{\lif}}
\BinaryInfC{$!C$}
\DisplayProof
\]
د دې پرېکون د لرې کولو دپاره په~$\delta$ کښې فرعي
ثبوت~$\delta_1$ په لاندې رسم بدلوو:
\[
\AxiomC{}
\RightLabel{$\delta_3$}
\DeduceC{$!B$}
\RightLabel{$\Subst{\delta_2}{\delta_3}{!B^x}$}
\DeduceC{$!C$}
\DisplayProof
\]
په $\Subst{\delta_2}{\delta_3}{!B^x}$ هغه !!{proof} مراد دے چې
له~$\delta_2$ څخه د هر~$!B$ فرض، چې نښه ئې~$x$ ده، د بشپړ
!!{proof}~$\delta_3$ په ځاے کولو حاصلېږي۔ په دې !!{proof} کښې
بدل شوي اصلي پرانيستي $!B$ فرضونه، چې نښه ئې~$x$ ده، نۀ دي
پاتې؛ خو دننه شوي ورته نښه لرونکي فرضونه پاتې کېدے شي۔ د~$x$
نښه لرونکے بل بېل فارمول نۀ شته، ځکه د معمول په شان فرض کوو
چې دوه بېل !!{formula}s د يوې نښې لرونکي فرضونه نۀ وي۔ د نوي
!!{proof} هر پرانيستے فرض يا له~$\delta_1$ څخه پاتې نۀ بدل شوے
پرانيستے فرض دے، يا له داخل شوي ثبوت څخه داسې فرض دے چې له
داخلېدو وروسته نۀ دے ايستل شوے۔ په~$\delta_2$ کښې د فرضونو
ايستونکے هر استنتاج په $\Subst{\delta_2}{\delta_3}{!B^x}$ کښې
د هماغې نښې فرضونه باسي؛ په هغو کښې نوي دننه شوي فرضونه هم
راتلے شي۔ په~$\delta_3$ کښې د فرضونو هر ايستونکے استنتاج ښايي
څو نقلونه ولري، چې په $\Subst{\delta_2}{\delta_3}{!B^x}$ کښې
اړوند فرضونه باسي۔
\textbf{د سرچينه‌يي دعوې سمون OLCMP-172، د OLCMP-133 تکراري مورد:}
د ټولو پخوانيو پرانيستو فرضونو عين پاتې کېدل لازم نۀ دي؛ دننه
شوے فرض ايستل کېدے شي، او که د ايستل کېدونکي فرض هېڅ مورد نۀ
وي، د آرګومېنټ ټول ثبوت حذفېږي۔ د پاتې فرضونو اصل دلته کره شو؛
د نښو د تازګۍ نوی شرط نۀ دے ورزيات شوے۔

هر پرېکون چې بشپړ په $\delta_2$ يا $\delta_3$ کښې وي، يا
په~$\delta$ کښې بشپړ له~$\delta_1$ څخه بهر وي، نۀ بدلېږي:
هماغه !!{formula}s رانغاړي او له هماغو استنتاجونو تېرېږي۔ نو
په ځانګړي ډول ئې رتبه او اوږدوالے په~$\delta$ کښې د اړوند
پرېکون سره عين دي۔ د تر ټولو لوړې رتبې يو د~$1$ اوږدوالي
پرېکون مو لرې کړے دے، نو يا د پرېکون اوږدوالے کم شوے يا ئې
رتبه ټيټه شوې ده، که د~$\delta$ د پرېکون اوږدوالے~$1$ و او
همدا د تر ټولو لوړې رتبې يوازينے پرېکون و۔
\textbf{د سرچينه‌يي حوالې سمون OLCMP-173:} دلته بېلګه يوازې
دويم او درېيم فرعي ثبوتونه لري؛ د څلورم ناتعريف فرعي ثبوت
کاپي شوې حواله لرې شوه۔

خو دوه خبرې کېدے شي:
\begin{enumerate}
  \item د ټولو $!B^x$ فرضونو په $\delta_3$ بدلولو سره مو
  په~$\delta_3$ کښې د ټولو پرېکونونو نقلونه زيات کړي دي۔ که په
  هغو کښې د تر ټولو لوړې رتبې پرېکونونه وي، د نوي !!{proof} د
  پرېکون اوږدوالے د زاړه !!{proof} له پرېکون اوږدوالي زياتېدے
  شي۔ بايد ډاډه شو چې داسې نۀ کېږي۔ دا په دې يقيني کوو چې د
  کمولو بدلونونه يوازې په تر ټولو \emph{ښي اړخ} د تر ټولو لوړې
  رتبې پرېکونونو تطبيق کړو: داسې پرېکونونه چې په~$\delta$ کښې
  ئې ښي اړخ ته په کومه څانګه د تر ټولو لوړې رتبې بل پرېکون نۀ
  وي۔ که په~$\delta_3$ کښې داسې پرېکون وي، د دې کمېدونکي
  پرېکون ښي اړخ ته به وي، نو دلته به مو تر ټولو ښي پرېکون نۀ
  و غوره کړے۔ تل د تر ټولو ښي پرېکون په غوره کولو سره يقيني
  کوو چې يا د !!{proof} د پرېکون اوږدوالے کم کړو يا ئې د پرېکون
  رتبه هم ټيټه کړو۔
  \item که په $\delta_2$ کښې فرض $!B^x$ د !!a{elimination}
  استنتاج لويه مقدمه وي او د~$\delta_3$ وروستے استنتاج
  \Elim{\lor}، \Elim{\lexists} يا !!a{introduction} وي، نو
  په~$\delta_2$ کښې پر داسې فرض د~$\delta_3$ په پېوندولو مو
  نوی پرېکون جوړ کړے دے۔ هغه څۀ چې په~$\delta_2$ کښې يوازې
  فرض و، اوس د~$1$ اوږدوالي پرېکون دے، يعنې د
  !!a{introduction} استنتاج پايله او د !!a{elimination} استنتاج
  لويه مقدمه، يا د پرېکون د برخې وروستے !!{formula} مورد دے،
  يعنې د !!a{elimination} استنتاج لويه مقدمه او د \Elim{\lor}
  يا \Elim{\lexists} پايله۔ که $B^x$ د \Elim{\lor} يا
  \Elim{\lexists} کوچنۍ مقدمه وه، له مخکې د يوې برخې، او ښايي
  د پرېکون د برخې، لومړے !!{formula} مورد و۔ په نوي !!{proof}
  کښې دا برخه اوس اوږدېدې شي۔ خو پام وکړئ چې د دې نوي پرېکون
  يا پخوانۍ پرېکون برخې !!{formula}~$!B$ دے او
  $\depth{!B} < \depth{!B \lif !C}$۔ نو که مو پرېکونونه زيات
  کړي يا ئې اوږدوالے زيات کړے هم وي، په هغو کښې يو هم د تر
  ټولو لوړې رتبې پرېکون نۀ دے۔ په دې توګه مو يا د پرېکون رتبه
  ټيټه کړې، يا که د هماغې رتبې نور پرېکونونه پاتې وي، د پرېکون
  اوږدوالے مو کم کړے دے۔
\end{enumerate}

د~$1$ اوږدوالي پرېکون د کمولو عمل د \emph{کمولو بدلون}
يا د \emph{تاو راتاو بدلون} بلل کېږي۔ د ځينو قاعده‌جوړو بڼې
په \olref{tab:reduction-conversions} کښې ورکړې دي؛ پاتې بڼې د
مشقونو په توګه پرېښودل شوې دي۔

\begin{table}
\begin{multline*}
\bottomAlignProof
\AxiomC{$\Discharge{!B}{x}$}
\RightLabel{$\delta_2$}
\shortDeduce
\DeduceC{$!C$}
\DischargeRule{\Intro{\lif}}{x}
\UnaryInfC{$!B \lif !C$}
\AxiomC{}
\RightLabel{$\delta_3$}
\shortDeduce
\DeduceC{$!B$}
\RightLabel{\Elim{\lif}}
\BinaryInfC{$!C$}
\DisplayProof
\quad \longrightarrow\quad
\shoveright{\bottomAlignProof
\AxiomC{}
\RightLabel{$\delta_3$}
\shortDeduce
\DeduceC{$!B$}
\RightLabel{$\Subst{\delta_2}{\delta_3}{!B^x}$}
\shortDeduce
\DeduceC{$!C$}
\DisplayProof}
\\[4ex]
\shoveleft{\bottomAlignProof
\AxiomC{}
\RightLabel{$\delta_2$}
\shortDeduce
\DeduceC{$!B$}
\AxiomC{}
\RightLabel{$\delta_3$}
\shortDeduce
\DeduceC{$!C$}
\RightLabel{\Intro{\land}}
\BinaryInfC{$!B \land !C$}
\RightLabel{\Elim{\land}[1]}
\UnaryInfC{$!B$}
\DisplayProof
\quad \longrightarrow\quad
\bottomAlignProof
\AxiomC{}
\RightLabel{$\delta_2$}
\shortDeduce
\DeduceC{$!B$}
\DisplayProof}\qquad
\shoveright{\bottomAlignProof
\AxiomC{}
\RightLabel{$\delta_2$}
\shortDeduce
\DeduceC{$!B$}
\AxiomC{}
\RightLabel{$\delta_3$}
\shortDeduce
\DeduceC{$!C$}
\RightLabel{\Intro{\land}}
\BinaryInfC{$!B \land !C$}
\RightLabel{\Elim{\land}[2]}
\UnaryInfC{$!C$}
\DisplayProof
\quad \longrightarrow\quad
\bottomAlignProof
\AxiomC{}
\RightLabel{$\delta_3$}
\shortDeduce
\DeduceC{$!C$}
\DisplayProof}
\\[4ex]
\bottomAlignProof
\AxiomC{}
\RightLabel{$\delta_2$}
\shortDeduce
\DeduceC{$!B$}
\RightLabel{\Intro{\lor}[1]}
\UnaryInfC{$!B \lor !C$}
\AxiomC{$\Discharge{!B}{x}$}
\RightLabel{$\delta_3$}
\shortDeduce
\DeduceC{$!D$}
\AxiomC{$\Discharge{!C}{y}$}
\RightLabel{$\delta_4$}
\shortDeduce
\DeduceC{$!D$}
\DischargeRule{\Elim{\lor}}{x\,y}
\TrinaryInfC{$!D$}
\DisplayProof
\quad \longrightarrow\quad
\shoveright{\bottomAlignProof
\AxiomC{}
\RightLabel{$\delta_2$}
\shortDeduce
\DeduceC{$!B$}
\RightLabel{$\Subst{\delta_3}{\delta_2}{!B^x}$}
\shortDeduce
\DeduceC{$!D$}
\DisplayProof}
\\[2ex]
\shoveleft{\bottomAlignProof
\AxiomC{}
\RightLabel{$\delta_2$}
\shortDeduce
\DeduceC{$!B(c)$}
\RightLabel{\Intro{\lforall}}
\UnaryInfC{$\lforall[x][!B(x)]$}
\RightLabel{\Elim{\lforall}}
\UnaryInfC{$!B(t)$}
\DisplayProof
\quad \longrightarrow\quad
\bottomAlignProof
\AxiomC{}
\RightLabel{$\Subst{\delta_2}{t}{c}$}
\shortDeduce
\DeduceC{$!B(t)$}
\DisplayProof}
\\[4ex]
\shoveright{\bottomAlignProof
\AxiomC{}
\RightLabel{$\delta_2$}
\shortDeduce
\DeduceC{$!B(t)$}
\RightLabel{\Intro{\lexists}}
\UnaryInfC{$\lexists[x][!B(x)]$}
\AxiomC{$\Discharge{!B(c)}{x}$}
\RightLabel{$\delta_3$}
\shortDeduce
\DeduceC{$!D$}
\DischargeRule{\Elim{\lexists}}{x}
\BinaryInfC{$!D$}
\DisplayProof
\quad \longrightarrow\quad
\bottomAlignProof
\AxiomC{}
\RightLabel{$\delta_2$}
\shortDeduce
\DeduceC{$!B(t)$}
\RightLabel{$\Subst{\Subst{\delta_3}{t}{c}}{\delta_2}{!B(t)^x}$}
\shortDeduce
\DeduceC{$!D$}
\DisplayProof}
\end{multline*}
\caption{د \Log{N1i} د کمولو ځينې بدلونونه۔}
\ollabel{tab:reduction-conversions}
\end{table}

\begin{prob}
د \Intro{\lnot} او ورپسې \Elim{\lnot} د کمولو بدلون ورکړئ۔
\end{prob}

يو حالت لا پاتې دے۔ د پرېکون تعريف هغه حالت هم رانغاړي
چې د پرېکون اوږدوالے~$1$ وي او $!A$ هم د~\FalseInt{} پايله
او هم د !!a{elimination} استنتاج لويه مقدمه وي۔ په دې حالت
کښې په هغه شان چلند کوو لکه چې $!A$ د !!a{introduction}
قاعدې پايله وي۔ مثلاً، دا رسم
\[
\bottomAlignProof
\AxiomC{}
\RightLabel{$\delta_2$}
\DeduceC{$\lfalse$}
\UnaryInfC{$!B \lif !C$}
\AxiomC{}
\RightLabel{$\delta_3$}
\DeduceC{$!B$}
\RightLabel{\Elim{\lif}}
\BinaryInfC{$!C$}
\DisplayProof
\quad\text{په دې}\quad
\bottomAlignProof
\AxiomC{}
\RightLabel{$\delta_2$}
\DeduceC{$\lfalse$}
\UnaryInfC{$!C$}
\DisplayProof
\]
په ښي اړخ رسم بدلوو۔ په هغو حالتونو کښې لږ احتياط پکار دے چې
$!A \ident (!B \lor !C)$ يا $!A \ident \lexists[x][!B(x)]$
وي۔ مثلاً، که لاندې چپ رسم په ښي رسم بدل کړو،
\[
\bottomAlignProof
\AxiomC{}
\RightLabel{$\delta_2$}
\DeduceC{$\lfalse$}
\RightLabel{\FalseInt}
\UnaryInfC{$!B \lor !C$}
\AxiomC{$\Discharge{!B}{x}$}
\RightLabel{$\delta_3$}
\DeduceC{$!D$}
\AxiomC{$\Discharge{!C}{y}$}
\RightLabel{$\delta_4$}
\DeduceC{$!D$}
\RightLabel{\Elim{\lor}}
\TrinaryInfC{$!D$}
\DisplayProof
\quad\text{په دې}\quad
\bottomAlignProof
\AxiomC{}
\RightLabel{$\delta_2$}
\DeduceC{$\lfalse$}
\RightLabel{\FalseInt}
\UnaryInfC{$!D$}
\DisplayProof
\]
نو د~$!D$ په پرېکون !!{formula} نوی پرېکون جوړېدے شي، او هېڅ
ضمانت نشته چې $\depth{!D} < \depth{!B \lor !C}$!{} نو دلته
هم د \Intro{\lor}/\Elim{\lor} د کمولو په شان عمل کوو او
په لاندې يوه بڼه ئې بدلوو:
\[
\bottomAlignProof
\AxiomC{}
\RightLabel{$\delta_2$}
\DeduceC{$\lfalse$}
\RightLabel{\FalseInt}
\UnaryInfC{$!B$}
\RightLabel{$\delta_3$}
\DeduceC{$!D$}
\DisplayProof
\quad\text{يا}\quad
\bottomAlignProof
\AxiomC{}
\RightLabel{$\delta_2$}
\DeduceC{$\lfalse$}
\RightLabel{\FalseInt}
\UnaryInfC{$!C$}
\RightLabel{$\delta_4$}
\DeduceC{$!D$}
\DisplayProof
\]
\textbf{د سرچينه‌يي نښې سمون OLCMP-174:} د يا ايستلو په
اصلي رسم کښې دواړو بېلو فرضونو يوه نښه لرله؛ د دويمې څانګې
نښه په بله نښه بدله شوه، څو د بېلو فارمولونو نښې سازګارې وي۔
نور ټول رسمي رسمونه عين ساتل شوي دي۔

\end{document}
